\documentclass[english,twoside, openright, hidelinks, 11 pt, class=report,crop=false]{standalone}
\input{../../preamb}
\input{../bm}

\begin{document}

\opgt

\op{rads}
Gjør om til radianer:\os
\abch{
\item $ 60^\circ $
\item $ 15^\circ $
\item $80^\circ$
\item $18^\circ$
}


\op{grad}
Gjør om til grader:\os
\abch{
\item $ \dfrac{11\pi}{12} $ 
\item $ \dfrac{11\pi}{6} $
\item $\dfrac{3\pi}{5}$
\item $\dfrac{\pi}{2}$
}

\eksop{R2V25}{r2v25d1opg7a}
$A$ er sentrum i sirkelen, som har radius 3. Buen mellom $B$ og $C$ har lengde 4.
\abc{
\item Finn verdien til $\angle BAC$ i radianer.
\item Finn verdien til $\angle BAC$ i grader.
}
\fig{r2v25d1opg7a}

\nes
\op{pytg}
Bruk Pytagoras' setning og definisjonen av $ \cos x $ og $ \sin x $ til å vise at
\[ \cos^2 x + \sin^2 x = 1 \] 

\op{tanx}
Finn $\tan x $ når du vet at\os
\abc{
\item $ \sin x = 0$ og $ \cos x=1 $ 
\item $ \sin x = \dfrac{1}{2} $ og $ \cos x = -\dfrac{\sqrt{3}}{2} $
}

\eksop{R2H24}{r2h24d1opg4}
Finn verdiene til $\cos x$ og $\tan x$ når ${\sin x=-\frac{1}{4}}$\\  og $x\in\left[\pi, \frac{3\pi}{2}\right]$.

\eksop{R2V25}{r2v25d1opg7b}
Finn verdiene til $\cos x$ og $\sin u$ når $\tan u=2$ \\ og $u\in\left(0, \frac{\pi}{2}\right)$.

\eksop{R2V26}{r2v26d1opg3}
Finn verdiene til $\sin v$ og $\tan v$ når $\cos v=\frac{2}{3}$, og $v$ er et tall i $4.$ kvadrant.



\newpage
\op{kvadrant}
Bruk $ - $ og $ + $ for å indikere henholdsvis negativ og positiv, og sett riktige markører i tabellen under.
\begin{center}
	\begin{tabular}{@{}l|l| l| l|l|}
		&1. kvadrant & 2. kvadrant & 3. kvadrant & 4. kvadrant\\ \hline
		$ \sin x $& & & &\\ \hline
		$ \cos x $& & & &\\ \hline
		$ \tan x $& & & &\\\hline
	\end{tabular}
\end{center}

\op{trigverd}
Bestem verdien til\os
\abch{
\item $ \sin\left(-\frac{\pi}{6}\right) $ 
\item $ \cos \pi $ 
\item $ \cos\left(-\frac{\pi}{2}\right) $
\item $ \tan\left(\frac{2\pi}{3}\right) $
}

\op{averdier}
Finn verdien til\os
\abch{
\item $ \asin 0 $ 
\item $ \asin \frac{\sqrt{3}}{2} $ 
\item $ \acos (-1) $ 
} \os
\abchs{4}{
\item $ \acos \left(-\frac{\sqrt{2}}{2}\right) $
\item $ \atan 1 $ 
\item $ \atan \left(\frac{1}{\sqrt{3}}\right) $
}


\op{bruk1}
Bruk \eqref{1} til å vise at $ {\sin\left(\frac{\pi}{6}\right)=\frac{1}{2}} $ når du vet at $ \cos\left(\frac{\pi}{6}\right)=\frac{\sqrt{3}}{2} $.

\op{sin2xopg} \vs
\abc{
\item Bruk én av de trigonometriske identitetene til å vise at
\[  \sin (2x)=2\cos x\sin x \]
\item Gitt at $ {\sin \left(\frac{\pi}{6}\right)=\frac{1}{2}} $. Bruk dette og identiteten over til å vise at
\[ 2\cos \left(\frac{\pi}{12}\right)\sin \left(\frac{\pi}{12}\right) = \frac{1}{2} \]
}


\op{cossomsin}
Skriv om utrykket
\[ \cos \left(3x-\frac{5\pi}{2}\right) \]
til et sinusuttrykk.
\newpage
\op{cossinsin}
Skriv om uttrykket \[ \cos( 2x) + \sqrt 3\sin (2x) \]
til et sinusuttrykk.

\nes

\op{sincos0}
Vis at alle løsninger av ligningen $ \cos x=0 $ er gitt som
\[ x = \frac{\pi}{2}+\pi n \]
mens alle løsninger av ligningen $ \sin x = 0 $ er gitt som
\[ x = \pi n \]\vs \vs

\op{triligns}
Løs ligningene:\os
\abc{
\item $ \cos x=\frac{\sqrt{2}}{2} $\

\item $ \cos \left(\frac{\pi}{3}x\right) = 0\quad,\quad x\in[0, 5]$

\item $ 2\sin(3x)=1 $\os

\item $ \sin(2x-\pi)=-\frac{\sqrt{3}}{2} $

\item $ 2\sqrt{3}\tan \left(4x+\frac{\pi}{2}\right) = 2$

\item $2\cos\left(\frac{\pi}{3}x\right)=\sqrt{3}\quad,\quad x\in(0, 10)$
}



\op{asinbcos0o}
Løs ligningene:
\abc{
\item $\sqrt{3} \sin x - \cos x =0$\os

\item $ \sin x + \sqrt{3}\cos x =0$\os

\item $ \cos (2x) + \sin (2x)=0  $\os
}




\op{kombo}
Løs ligningene:\os
\abc{
\item $ \cos x - \sin x = \sqrt{2}  $ 
\item $ \sqrt{3} \cos\left(\dfrac{x}{2\pi}\right) - \sin\left(\dfrac{x}{2\pi}\right)=1 $
\item $\sin x-\sqrt{3}\cos x=0\qquad,\qquad x\in[0,2\pi)$
}

\nes
\newpage
\op{abctri}
Løs ligningene:\\
\abc{
\item $ \sin^2 x+\frac{1}{2}\sin x-\frac{1}{2}=0 $

\item $ \cos^2 (3x) -3\cos (3x) -4=0$

\item $ 2\cos^2 x+ \sqrt{8} \cos x+1=0 $

\item $ \tan^2 (\pi x) -\sqrt{12}\tan (\pi x) +  3=0$

\item $ -\sin^2(3x)- 3\cos (3x) -3=0$
}



\op{kvadopg} 
Løs ligningene:\os
\abc{
\item $ -\cos^2 x+15\sin^2 x = 3 $ 
\item $ \cos^2 \left(\dfrac{x}{4}\right) - 2 \sin^2 \left(\dfrac{x}{4}\right) = -\dfrac{1}{2} $
}


\eksop{R2V23}{r2v23d1opg3}
Løs likningen
\[ \cos^2 x-3\sin^2 x=-2 \quad,\quad x\in[0, 2\pi)\]

\nes

\op{cosmaks}	
Forklar hvorfor en cosinusfunksjon $ f(x)=a\cos(kx+c)+d $ har\os
\abc{
\item maksimalverdi for $ kx+c= 2\pi n$ og minimalverdi for $ kx+c= \pi + 2\pi n $ når $ a>0 $.
\item maksimalverdi for $ kx+c= \pi+2\pi n$ og minimalverdi for $ kx+c= 2\pi n $ når $ a<0 $.
\item ekstremalverdi for $kx+c=\pi n$.
}



\op{cosfunko}
Gitt funksjonen 
\[ f(x)=-3\cos\left(3x+\frac{\pi}{12}\right)+4 \]
\abc{
\item Finn perioden til $ f $.

\item Hva er minimums- og maksimumsverdiene til $ f $?

\item Finn alle $ x $ hvor $ f $ har minimums- og maksimumsverdier.
}


\op{opgtrifinnmaks}
Finn maksimalpunktene til funksjonen $ f(x)= 2 \cos (3x+\pi)+1 $.

\newpage
\op{sinmaks}
Forklar hvorfor en sinusfunksjon $ f(x)=a\sin(kx+c)+d $ har\os
\abc{
\item maksimalverdi for $ kx+c=\frac{\pi}{2}+ 2\pi n$ og minimalverdi for $ kx+c= -\dfrac{\pi}{2} + 2\pi n $ når $ a>0 $.
\item maksimalverdi for $ kx+c= -\frac{\pi}{2} + 2\pi n$ og minimalverdi for $ kx+c=\dfrac{\pi}{2}+ 2\pi n $ når $ a<0 $. 
\item ekstremalverdi for $kx+c=\frac{\pi}{2}+\pi n$.
}



\op{sinfnpkt}
Gitt funksjonen 
\[ f(x)=-2\sin\left(\frac{\pi}{2}x\right)+1 \quad, \quad x\in[-5, 5]\]
\abc{
\item Finn perioden til $ f $.

\item Finn topppunktene til $ f $.

\item Finn nullpunktene til $ f $.
}


\op{skissin}
Om en cosinusfunksjon vet du følgende:
\begin{itemize}
	\item likevektslinja til funksjonen er $ y=1 $.
	\item $ (0,3) $ og $ (2\pi, 3) $ er to naboliggende toppunkt.
\end{itemize}
Skisser grafen til funksjonen for ${ x\in[0, 3\pi]} $.

\op{opgtriomskriv1} \vs
\abc{
\item Gitt funksjonen
\[ f(x)=a\cos(kx+c)+d \]
Vis at vi kan skrive
\[ f(x)=a\sin\left(kx+c+\frac{\pi}{2}\right)+d \]
\item Gitt funksjonen
\[ g(x)=a\sin(kx+c) \]
Vis at vi kan skrive
\[ g(x)=a\cos\left(kx+c-\frac{\pi}{2}\right) \]
}




\op{opgtriomskriv}
Skriv om funksjonen $ f(x)=2 \cos (3x+\pi)+1 $ til en sinusfunksjon.

\newpage

\op{sincosfig}
Figuren under viser grafen til funksjonen $f(x)$.
\abc{
\item Uttrykk $f$ som en cosinusfunksjon.
\item Uttrykk $f$ som en sinusfunksjon.
}

\fig{cosopg2}

\eksop{R2V26}{r2v26d1opg4}
Figuren under viser grafen til en trigonometrisk funksjon $f$. Finn et uttrykk for $f$.
\fig{r2v26d1opg4}

\op{tanfopg}
Forklar hvorfor alle funksjoner $ f $ gitt ved
\[ f(x)= a \tan (kx+c)+d\]
har vertikale asymptoter når
\[ kx+c=\pm \frac{\pi}{2}+\pi n \]


\eksop{R2V24}{r2v24d1opg5}
Gitt funksjonen
\[ f(x)=2\sin\left(\frac{\pi}{6}x-\frac{\pi}{3}\right)-1\quad,\quad D_f=(0, 20) \]
\abc{
\item Løs likningen $f(x)=0$.
\item Finn amplituden, likevektslinja, perioden og forskyvningen langs likevektslinja.
}

\eksop{R2V25}{r2v25d1opg4}
Gitt funksjonen
\[ f(x)=2\sqrt{3}\sin\left(2x+\frac{\pi}{6}\right)\qquad,\qquad D_f=\left(0, \frac{\pi}{2}\right) \]
\abc{
\item Finn amplituden, likevektslinja, perioden og faseforskyvningen til $f$.
\item Løs likningen $f(x)=\sqrt{3}$.
}
Gitt funksjonen
\[ g(x)=3\sin(2x)+\sqrt{3}\cos(2x)\qquad,\qquad D_g=\left(0, 2\pi\right) \]
\abcs{3}{
\item Løs likningen $g(x)=\sqrt{3}$.	
}

\op{opgtriderfunk}
Finn den deriverte av funksjonen $ f(x)=\frac{\cos x}{x^4} $.
\newpage

\grubop{udir}
En annen vanlig definisjon av radianer er følgende:
\st{
Vinkel $ v $ mellom to linjestykker $ AB $ og $ AC $, målt i radianer, tilsvarer forholdet mellom lengden på en bue mellom $ AB $ og $ AC $ og radien til buen.
}
I figuren under er det tegnet en bue mellom $AB$ og $AC$. Vi setter lengden og radien til buen henholdsvis lik $l$ og $r$. I følge definisjonen over er altså $ v=\dfrac{l}{r} $.
\os
Forklar hvorfor radianer ut ifra denne definisjonen også kan sees på som en buelengde langs enhetssirkelen.

\begin{figure}
	\centering
	\includegraphics[]{\figp{rad}}
\end{figure}




\grubeksoprecc{r2h24d1opg5}{R2H24}
Figuren viser grafen til funksjonen
\[ f(x)=2\sin\left(\frac{\pi}{4}x-\frac{\pi}{2}\right)-1 \]
\abc{
\item Bestem en funksjon $g(x)=A\cos(cx+\phi)+d$ som passer til grafen.
\item Løs likningen $\sin\left(\frac{\pi}{4}x-\frac{\pi}{2}\right)=\frac{1}{2}$, der $x\in\left[0, 3\pi\right]$. \\ Hvilken type punkt er disse $x$-verdiene for funksjonen $f$?
}
\fig{r2h24d1opg5}

\grubop{eksmpr2v23d2opg2}
En uendelig geometrisk rekke er gitt ved
\[ S(x)=\cos x+2\sin x\cos x + 4\sin^2 x\cos x+... \]
\abc{
\item Bestem området i intervallet $[0, \pi]$ hvor rekken konvergerer.
\item Avgjør for hvilke verdier av $b$ likningen 
\[ S(x)=b \]
har én eller flere løsninger i intervallet $[0, \pi]$ 
}


\grubeksoprecc{r2h25d1opg4}{R2H25}
Gitt likningen
\[ \left(\sin x-\frac{1}{2}\right)(\sin x-a)=0\qquad,\qquad x\in[0, 2\pi) \]
Hvor hvilke verdier av $a$ har likningen henoldsvis to, tre og fire løsninger?


\grubeksop{r2v23d1opg5}{R2V23D1}
I denne oppgaven skal du vise at $ \lim\limits_{v\to 0^+} \dfrac{\sin v}{v}=1 $.\os

I figuren nedenfor er $ AB=AD=1 $, og buen mellom $ B $ og $ D $ er del av en sirkel med sentrum i $ A $. Vi lar $ \angle BAC=v $.
\fig{r2v23d1opg5}
\abc{
	\item Bruk arealbetraktninger til å grunngi at
	\[ \frac{1}{2}\sin v<\frac{1}{2}v<\frac{1}{2}\tan v \]
	\item Forklar at dette gir oss
	\[ 1<\frac{v}{\sin v}<\frac{1}{\cos v} \]
	\item bruk ulikhetene fra oppgave b) til å grunngi at $ \lim\limits_{v\to 0^+} \dfrac{\sin v}{v}=1 $.
}


\grubop{opgtangr}
Gitt at
\[ \tan x = \frac{a}{b} \]
Vis at $ \sin x = \dfrac{a}{\sqrt{a^2+b^2}} $ og at $ \cos x = \dfrac{b}{\sqrt{a^2+b^2}} $.

\grubop{opgtridet}
Gitt vektorene $ \vec{u}=[a, b] $ og $ \vec{v}=[c, d] $. Vis at
\[ |\det(\vec{u}, \vec{v})|=|u||v|\sin\angle(\vec{u}, \vec{v}) \]

\begin{comment}

\ekspop
I denne oppgaven skal vi komme fram til to viktige resultater, nemlig at:
\alg{
\lim\limits_{x\to 0}\frac{\sin x}{x}&=1 \\ &\\
\lim\limits_{x\to 0}\frac{\cos x-1}{x}&=0 	
}
Når vi i \textsl{Kapittel 6} finner den deriverte av $ \sin x $ og $ \cos x $, brukes blant annet disse resultatene.
\begin{figure}[H]
\centering
\includegraphics[]{\asym{sinx}}
\end{figure}
a) Bruk formlikhet til å vise at $ \tan x $ er høyden i figuren over.

b) Forklar hvorfor vi må ha at:
\[ \frac{\sin x}{x}<1 \]
c) Buen $ x $ utgjør $ \frac{x}{2\pi} $ av omkretsen til enhetssirkelen. Vis at arealet av sektoren til $ x $ blir $ \frac{1}{2}x $.

d) Forklar hvorfor vi må ha at:
\alg{
\frac{1}{2}x&<\frac{1}{2}\tan x	\\
x&<\tan x	
}
e) Bruk ulikhetene fra \textsl{b} og \textsl{d} til å vise at:
\[ \lim\limits_{x\to 0} \frac{\sin x}{x} =1 \]

f) Vis at:
\[\lim\limits_{x\to 0} \frac{\cos x-1}{x}=0 \]
\big(Hint: Multipliser ligningen med $ \frac{\cos x+1}{\cos x+1} $ og bruk deretter det du fant i $ e $.\big) \\
\end{comment}

\end{document}